\section{A directed chain of conditional CRFs}
\label{sec:construction}

\TODO{draft from \texttt{causalprobabilistictransformer\_1.pdf}, Part~II.}

\subsection{The chain}

Instead of one CRF over the sequence, the model is a chain of conditional CRFs, one per
position. At step $t$ the CRF over $(\Zv{t}, \Hv{t}{1..\nchan}, \Wv{t})$ is conditioned on the
prefix posteriors $\{\QZ{s}\}_{s < t}$, which enter as fixed numbers:
%
\begin{equation}
\label{eq:chain}
  \TODO{the factorisation, with $\stopgrad{\cdot}$ marking the frozen prefix}
\end{equation}

\subsection{The prefix as a frozen condition}
\label{sec:frozen-prefix}

\FROMPDF{Part~II --- the precise conditioning statement}
Each $\QZ{s}$, $s<t$, is treated as data at step $t$: no message travels from $t$ back to $s$,
and no gradient does either. Three things follow.

\begin{itemize}
  \item The anti-causal term of \Cref{thm:trilemma} is never formed
        (\Cref{rmk:never-formed}), so \Cref{def:causality} holds by construction rather than by
        masking.
  \item The per-step partition function is over step $t$'s variables only, which is what makes
        the exact readout of \Cref{sec:exact-readout} available.
  \item Training is still a single backward pass over the sequence: the frozen condition
        removes a gradient \emph{path}, not the gradient.
\end{itemize}

\TODO{state explicitly what is lost. Being honest here is cheaper than being caught: the chain
is no longer a single joint distribution over the sequence, so quantities defined by that joint
--- e.g.\ an exact sequence-level marginal --- are not available. Argue that this is the same
trade a standard autoregressive factorisation makes.}

\subsection{Encoding and decoding as two conditioning patterns}
\label{sec:one-model}

\TODO{state here, prove in \Cref{sec:output}: the encoder clamps $\Wv{t}$ and leaves the
sequence undirected; the decoder unclamps $\Wv{t}$ and directs the chain. Same parameters, same
factors, two conditioning patterns --- not two models. This is what constraint~4 of the project
means and it should be visible in the paper, since it is the cleanest statement of the
contribution.}
